% Fichier généré automatiquement par tools/generate_corrections.py.
% Source : STAT/STAT-03-Demarche/55_Suspension/55_Suspension.tex
% Statut : complet (2/2 question(s) renseignée(s), 0 reprise(s)).
\begin{corrigebox}[Corrigé extrait de la banque]
\CorrigeQuestion{1}
\par\medskip\begin{center}\captionsetup{type=figure}
\includegraphics[width=10cm]{55_01_cor}
\end{center}\par\medskip
\CorrigeQuestion{2}
Calculons le degré d'hyperstatisme : 
\begin{itemize}
\item mobilités : $m=2$ (rotations autour de $\vect{a}$ et de $\vect{z}$);
\item inconnues statiques : $I_s = 3 \times 4 = 12$;
\item équations : $E_s = 2\times 6=12$. 
\item $h = m-E_s + I_s= 2 -12 + 12 = 2$.
\end{itemize}
On ne peut donc pas déterminer toutes les actions mécaniques.
\end{corrigebox}
